\documentclass[11pt,a4paper]{article}
\usepackage[utf8]{inputenc}
\usepackage[T1]{fontenc}
\usepackage[margin=2cm]{geometry}
\usepackage{enumitem}
\usepackage{xcolor}
\usepackage{hyperref}

% Colors
\definecolor{accent}{HTML}{2E86AB}
\definecolor{heading}{HTML}{000000}

% Section formatting
\usepackage{titlesec}
\titleformat{\section}
  {\Large\bfseries\color{heading}}
  {}{0em}{}[\titlerule]

% Better line spacing
\usepackage{setspace}
\setstretch{1.1}

% Custom commands with proper indentation
\newcommand{\cvitem}[3]{%
  \vspace{0.3em}%
  \noindent\textbf{#1} \hfill #2 \\
  \noindent\textit{#3}%
  \vspace{0.2em}%
}

\newcommand{\cvdetail}[1]{%
  \begin{adjustwidth}{1em}{0em}
    #1
  \end{adjustwidth}%
  \vspace{0.5em}%
}

\newcommand{\cvsubitem}[2]{%
  \vspace{0.3em}%
  \noindent\textbf{#1} \hfill #2%
  \vspace{0.2em}%
}

% For the adjustwidth environment
\usepackage{changepage}

% Better itemize spacing
\setlist[itemize]{leftmargin=*, itemsep=0.1em, parsep=0.1em, topsep=0.2em}

\begin{document}

% Header
\begin{center}
  {\Huge\bfseries Gerardo Cicalese}\\[0.3em]
  {\large Ph.D. Student in Mathematical Models and Methods in Engineering}\\[0.2em]
  Politecnico di Milano \quad | \quad 
  \href{mailto:gerardo.cicalese7@gmail.com}{gerardo.cicalese7@gmail.com}
\end{center}

\vspace{0.5em}

% Profile
\section{Profile}
Ph.D. student in Mathematical Models and Methods in Engineering with research experience in numerical simulation of wave propagation, partial differential equations, and inverse problems. Experienced in the development, implementation, and validation of numerical methods and scientific simulation software using MATLAB, Python, and C/C++. Research activities include finite-difference and finite-element methods, spectral methods, domain decomposition, waveform relaxation, differentiable physics, and parallel computing.

% Technical Expertise
\section{Technical Expertise}
\begin{itemize}
  \item \textbf{Programming:} C++, C, MATLAB, Python, \LaTeX
  \item \textbf{Numerical Methods:} FEM, FDTD, PSTD, spectral methods, domain decomposition, waveform relaxation, differentiable physics
  \item \textbf{Scientific Computing:} OpenMP, CUDA/OpenACC, Intel VTune Profiler, COMSOL Multiphysics
  \item \textbf{Analysis:} PDEs, numerical stability, convergence analysis, inverse problems, performance analysis
  \item \textbf{Systems:} Linux/Unix, Windows
  \item \textbf{Languages:} Italian (native), English (C1 Cambridge Certified), French (basic comprehension)
\end{itemize}

% Research Experience
\section{Research Experience}
\cvsubitem{Ph.D. Researcher -- Numerical Simulation and Scientific Computing}{2023--present}
\cvdetail{Politecnico di Milano -- MOX, Modeling and Scientific Computing Laboratory}
\begin{itemize}
  \item Developed and analyzed numerical methods for the simulation of wave propagation governed by partial differential equations.
  \item Implemented and validated numerical solvers in MATLAB and C/C++, including finite-difference and spectral methods.
  \item Developed domain decomposition and waveform relaxation techniques for large-scale wave simulations, including numerical and asymptotic optimization of transmission conditions.
  \item Investigated inverse problems for parameter reconstruction through Differentiable Physics Neural Networks (DPNNs).
  \item Investigated numerical stability, convergence, accuracy, and computational efficiency through systematic numerical experiments.
  \item Developed scientific simulation software and reproducible computational workflows for research applications.
  \item Applied parallel computing techniques using OpenMP and analyzed computational performance with Intel VTune Profiler.
\end{itemize}

\newpage
% Education
\section{Education}

\cvitem{Ph.D. in Mathematical Models and Methods in Engineering}{2023--present}
{Politecnico di Milano, Italy}
\cvdetail{Thesis: Domain Decomposition and Differentiable Physics for the Viscoelastic Wave Equation.\\
Supervisors: Prof. Ilario Mazzieri, Prof. Gabriele Ciaramella\\
Research Group: MOX - Modeling and Scientific Computing Laboratory}

\cvitem{M.Sc. in Music and Acoustic Engineering}{2021--2023}
{Politecnico di Milano, Italy}
\cvdetail{Thesis: Efficient techniques for accurate sound propagation modeling in virtual acoustics.\\Final mark: 110L/110.}

\cvitem{B.Sc. in Computer Engineering}{2018--2021}
{University of Salerno, Italy}
\cvdetail{Thesis: Analisi e test di Jukebox, una piattaforma per la generazione di pattern musicali.\\Final mark: 110L/110.}

% Selected Projects
\section{Selected Projects}

\cvsubitem{Personal GitHub Repository}{}
\cvdetail{\url{https://github.com/mazamin7} - Public code repositories including:}

\begin{itemize}
  \item \textbf{WASAbi -- Wave-based Acoustic Simulator using ARD} 
        \newline C++ scientific computing project implementing Adaptive Rectangular Decomposition for numerical simulation of acoustic wave propagation. Includes frequency-dependent atmospheric absorption, domain decomposition, numerical time integration, and reproducible simulation workflows. The implementation includes numerical stability and convergence studies and comparison of alternative spatial discretizations.
        
  \item \textbf{Damped Wave Equation SWR}
        \newline MATLAB implementation of Schwarz Waveform Relaxation (SWR) for the viscoelastic-telegrapher damped wave equation. Features spectral parameter optimization ($L^2$ and $L^\infty$) and reproducibility scripts for research presented at the DD29 conference.

  \item \textbf{SonicBlend}
        \newline Cross-synthesis application in MATLAB for audio signal processing. Implements LPC-based cross-synthesis with configurable parameters for carrier/modulator analysis and multiple solution methods including Wiener-Hopf and gradient descent.
\end{itemize}

% Selected Publications
\section{Selected Publications}

\cvsubitem{Addressing Atmospheric Absorption in Adaptive Rectangular Decomposition}{2024}
\cvdetail{\textit{The Journal of the Acoustical Society of America}, 156(4):2328--2339.\\Gerardo Cicalese, Gabriele Ciaramella, Ilario Mazzieri.}

\cvsubitem{Optimized Schwarz Waveform Relaxation for the Damped Wave Equation}{2026}
\cvdetail{\textit{Proceedings of the 29th International Conference on Domain Decomposition (DD29)} / arXiv preprint.\\Gerardo Cicalese, Gabriele Ciaramella, Ilario Mazzieri, Martin J. Gander.}

\cvsubitem{Asymptotic Optimization of Transmission Conditions for Schwarz Waveform Relaxation Applied to the Damped Wave Equation}{In prep.}
\cvdetail{Gerardo Cicalese, Gabriele Ciaramella, Ilario Mazzieri, Martin J. Gander.}

\cvsubitem{Sound Field Reconstruction through Differentiable Physics}{In prep.}
\cvdetail{Alinda Ezgi Gerçek, Federico Miotello, Gerardo Cicalese, Mirco Pezzoli, Fabio Antonacci, Ilario Mazzieri.}

\newpage
% Additional Experience
\section{Additional Experience}

\cvsubitem{Conferences and Presentations}{}
\begin{itemize}
  \item \textbf{WCCM--ECCOMAS 2026 (Munich):} Talk on Asymptotic Optimization of Transmission Conditions for Schwarz Waveform Relaxation Applied to the Damped Wave Equation.
  \item \textbf{Swiss Numerical Analysis Day 2025 (Basel):} Poster on Enhancing ARD with Atmospheric Absorption Modeling.
  \item \textbf{Domain Decomposition 29 (DD29) (Milan):} Poster on ARD Atmospheric Absorption Modeling and Talk on Optimized Red-Black Waveform Relaxation for Damped Wave Equation.
  \item \textbf{SIMAI 2025 (Trieste):} Talk on Optimized Red-Black Waveform Relaxation for Damped Wave Equation.
\end{itemize}

\cvsubitem{Research Stays}{}
\begin{itemize}
    \item \textbf{Université de Genève (Switzerland, 2025):} Research visits supervised by Prof. Martin J. Gander (March--May 2025; September--November 2025).
\end{itemize}

\cvsubitem{Teaching Experience}{}
\begin{itemize}
    \item \textbf{Teaching Assistant:} Curve e superfici per il design (Politecnico di Milano, 2023--2026).
    \item \textbf{Teaching Assistant:} Calcolo numerico (Politecnico di Milano, 2023--2024).
\end{itemize}

\cvsubitem{Professional Development \& Memberships}{}
\begin{itemize}
    \item \textbf{PhD Coursework:} Advanced Numerical Methods for Predictive Digital Twins; Mathematical and Numerical Foundations of Scientific Machine Learning; Parallel Computing through OpenMP and OpenACC / CUDA.
    \item \textbf{CIME School "PDEs, Control and Deep Learning" (July 2024):} Intensive course on PDEs and machine learning.
    \item \textbf{Domain Decomposition Conference Courses (June 2025):} Pre-conference training.
    \item \textbf{SIMAI:} Società Italiana di Matematica Applicata e Industriale.
\end{itemize}

\cvsubitem{Additional Education}{}
\begin{itemize}
    \item \textbf{Scientific High School Diploma:} Liceo Scientifico "Nicola Sensale" (2013--2018). Final mark: 98/100.
\end{itemize}

\cvsubitem{Other Software Projects}{}
\begin{itemize}
  \item \textbf{Rhythm Wheel, Chordophone Champion, FM Synth, YetAnotherAutoWah:} Various audio, synthesis, and interaction design projects using Vue.js, Tone.js, SuperCollider, and JUCE.
\end{itemize}

\cvsubitem{Media Coverage}{}
\begin{itemize}
    \item \textbf{La Provincia di Cremona (2023):} \href{https://www.laprovinciacr.it/news/cronaca/418996/lacustica-per-il-metaverso-gli-ingegneri-protagonisti.html}{"L'acustica per il metaverso: gli ingegneri protagonisti"}
    \item \textbf{Insieme News (2023):} \href{https://www.insiemenews.it/2023/10/23/inseguendo-il-suono/}{"Inseguendo il suono"}
\end{itemize}

\end{document}